% TOP_PROBLEM: {"id": 30004201, "title": "Chromatic Fixed-Point Implications for Finite p-Groups", "category_id": 7, "category": {"id": 7, "name": "topology", "display_name": "Topology", "description": "Properties preserved under continuous deformations.", "slug": "topology", "order_index": 7, "created_at": "2026-07-31T15:26:25.671Z"}, "status": "partially_solved"}
% FIRST_POSTED: 2026-10-06
% ATTEMPT_STATUS: solved
% OPENAI_CLAIM: {"schema_version":1,"source_repo":"https://github.com/openai/math","source_commit":"adc7f1241b42e322a6451854ab7e4b4c146bf78a","release_date":"2026-10-06","match":"full","resolution":"proved","independently_reviewed":false,"second_pass":"agreed","families":[{"family":"314","title":"Cyclic length and chromatic fixed-point loss","subject":"Topology","manuscripts":[{"title":"Cyclic length and chromatic fixed-point loss","pdf_path":"preprints/Cyclic-Length-and-Chromatic-Fixed-Point-Loss-September-24-2026/Cyclic-Length-and-Chromatic-Fixed-Point-Loss-September-24-2026.pdf","readme_path":"preprints/Cyclic-Length-and-Chromatic-Fixed-Point-Loss-September-24-2026/README.md","date":"2026-09-24","theorem_ref":"thm:main (sections/01-introduction.tex:18)","statement_sources":[{"path":"preprints/Cyclic-Length-and-Chromatic-Fixed-Point-Loss-September-24-2026/build/sections/01-introduction.tex","line":18,"label":"thm:main"}]}],"lean":null,"reasoning_trace":null}]}
\documentclass[11pt]{article}
\usepackage[margin=25mm]{geometry}
\usepackage{fontspec}
\setmainfont{Latin Modern Roman}
\usepackage{amsmath,amssymb,mathtools}
\usepackage{unicode-math}
\setmathfont{Latin Modern Math}
\usepackage{enumitem}
\usepackage{xurl,hyperref,bookmark}
\hypersetup{colorlinks=true,urlcolor=blue}
\setcounter{secnumdepth}{0}
\setlength{\parindent}{0pt}
\setlength{\parskip}{5pt}
\setlength{\emergencystretch}{3em}
\begin{document}
\section*{Chromatic Fixed-Point Implications for Finite p-Groups}

\subsection{OpenAI's claim}
For every subgroup of a finite p-group, the optimal chromatic loss under fixed points is claimed to equal its shortest subnormal-chain length with cyclic quotients, independently of height. The cited manuscript argues for that result; its relation to this problem is specified below.

\subsection{Mathematical statement of the claimed result}
\paragraph{Mathematical claim 1.} Fix a prime $p$, finite $p$-group $G$, subgroup $H$ and integer $n\ge0$. Let $r_n(G,H)$ be the least $r\ge0$ such that for every finite $p$-local genuine $G$-spectrum $X$,
\[
K(n+r)_*(\Phi^HX)=0\quad\Longrightarrow\quad K(n)_*(\Phi^GX)=0.
\]
Here $\Phi$ denotes geometric fixed points and $K(0)=H\mathbb Q$. Let $\ell(G,H)$ be the smallest number of steps among subnormal chains from $H$ to $G$ having a cyclic quotient at each step, and $\ell(G,G)=0$. OpenAI claims $r_n(G,H)=\ell(G,H)$ for all these parameters. Minimality includes finite witness spectra ruling out every smaller loss.

\subsection{Outline of the argument}
Manuscript 1 (family 314, 2026-09-24): Character evaluations construct a pro-p comparison group whose finite sections need few generators. Groupoid integration supplies subgroup-count congruences that force a sufficiently short closed series with Z\_p factors. Finite chromatic fracture translates this arithmetic bound into fixed-point vanishing, and tensor-prime separation provides witnesses showing that no smaller loss can work.

\subsection{Scope relative to this problem}
Manuscript 1 (family 314, 2026-09-24): The sharp loss r\_n(G,H)=ell(G,H) gives precisely m>=n+ell(G,H) for the finite-spectrum acyclicity implication; type monotonicity handles larger m, and trivial-action type-n spectra disprove m<n.

\subsection{Formal verification}
Family 314 has no Lean scope page in this snapshot; no formalization of these manuscript results is listed.

This site has not executed the Lean files or checked their dependencies. These descriptions report the scope of the linked declarations, rather than a new verification of the proofs.

\subsection{Status caveat}
OpenAI attributes these results to an internal model. Its release notice cautions that results lacking full formalization can contain mistakes. This site has not independently assessed the proofs; a solved label records the scope of the reported claim.

\subsection{Sources}
\begin{itemize}
\item OpenAI, \emph{Cyclic length and chromatic fixed-point loss}, 2026-09-24; citation key \texttt{OAI:Cyclic-Length-and-Chromatic-Fixed-Point-Loss-September-24-2026}. \href{https://github.com/openai/math/blob/adc7f1241b42e322a6451854ab7e4b4c146bf78a/preprints/Cyclic-Length-and-Chromatic-Fixed-Point-Loss-September-24-2026/Cyclic-Length-and-Chromatic-Fixed-Point-Loss-September-24-2026.pdf}{PDF (GitHub)}; \href{https://github.com/openai/math/raw/adc7f1241b42e322a6451854ab7e4b4c146bf78a/preprints/Cyclic-Length-and-Chromatic-Fixed-Point-Loss-September-24-2026/Cyclic-Length-and-Chromatic-Fixed-Point-Loss-September-24-2026.pdf}{PDF (direct)}; \href{https://github.com/openai/math/blob/adc7f1241b42e322a6451854ab7e4b4c146bf78a/preprints/Cyclic-Length-and-Chromatic-Fixed-Point-Loss-September-24-2026/README.md}{Manuscript page and citation}; \href{https://github.com/openai/math/blob/main/preprints/Cyclic-Length-and-Chromatic-Fixed-Point-Loss-September-24-2026/Cyclic-Length-and-Chromatic-Fixed-Point-Loss-September-24-2026.pdf}{Latest version}; Statement source: \href{https://github.com/openai/math/blob/adc7f1241b42e322a6451854ab7e4b4c146bf78a/preprints/Cyclic-Length-and-Chromatic-Fixed-Point-Loss-September-24-2026/build/sections/01-introduction.tex\#L18}{thm:main}.
\item \href{https://github.com/openai/math/blob/adc7f1241b42e322a6451854ab7e4b4c146bf78a/CONTENTS.md}{OpenAI catalogue entry}.
\end{itemize}
\end{document}
